\section{Late-Stage Real-Data Results and Negative Ablations}\label{app:late-results}

This appendix preserves the detailed late-stage tables that inform the main-text claim audit. These experiments are not presented as independent external validations. The locked TCGA-BRCA/METABRIC experiment preceded the cross-cohort pathway, anti-underclustering, and teacher--student development analyses.

\subsection{End-to-end stability-regularized transport}
\begin{table*}[t]
\centering
\caption{End-to-end stability-regularized UOT on the five real datasets. Values are mean $\pm$ standard deviation over five seeds. Counts were inherited from the label-free V4.1 consensus.}
\label{tab:v42-appendix}
\resizebox{\textwidth}{!}{%
\begin{tabular}{lcccccc}
\toprule
Dataset & NMI & ARI & ACC & Macro-F1 & Rare recall & $\widehat K/K$ \\
\midrule
TCGA-PANCAN RNA-Seq & 0.924 $\pm$ 0.025 & 0.901 $\pm$ 0.062 & 0.964 & 0.969 & 0.956 & 5/5 \\
Mice Protein & 0.422 $\pm$ 0.014 & 0.230 $\pm$ 0.012 & 0.343 & 0.214 & 0.000 & 4/8 \\
WDBC & 0.570 $\pm$ 0.147 & 0.667 $\pm$ 0.137 & 0.907 & 0.899 & 0.829 & 2/2 \\
Dermatology & 0.834 $\pm$ 0.014 & 0.766 $\pm$ 0.018 & 0.825 & 0.747 & 0.870 & 5/6 \\
Parkinsons & 0.103 $\pm$ 0.018 & 0.054 $\pm$ 0.026 & 0.800 & 0.692 & 0.429 & 6/2 \\
\bottomrule
\end{tabular}}
\end{table*}


The full and no-stability variants had nearly identical agreement metrics across 25 matched runs. Mean temporal stability increased from 0.823 to 0.832 ($p=0.0318$), whereas NMI ($p=0.776$) and ARI ($p=0.955$) did not improve.

\subsection{Pathway-representation benchmark}
\begin{table*}[t]
\centering
\caption{Label-free consensus results for the non-PAM50 gene and pathway representations. These experiments used both cohorts for development and should not be interpreted as external validation.}
\label{tab:v50-appendix}
\resizebox{\textwidth}{!}{%
\begin{tabular}{llccccc}
\toprule
View & Cohort & $\widehat K$ & NMI & ARI & Rare recall & Partition stability \\
\midrule
Gene PCA & TCGA-BRCA & 5 & 0.310 & 0.243 & 0.194 & 0.230 \\
Gene PCA & METABRIC & 5 & 0.085 & 0.067 & 0.100 & 0.126 \\
Hallmark & TCGA-BRCA & 2 & 0.034 & 0.056 & 0.000 & 0.822 \\
Hallmark & METABRIC & 2 & 0.091 & 0.086 & 0.000 & 0.276 \\
PROGENy & TCGA-BRCA & 3 & 0.208 & 0.257 & 0.000 & 0.580 \\
PROGENy & METABRIC & 3 & 0.196 & 0.231 & 0.000 & 0.661 \\
Pathway fused & TCGA-BRCA & 3 & 0.202 & 0.221 & 0.000 & 0.386 \\
Pathway fused & METABRIC & 3 & 0.235 & 0.222 & 0.000 & 0.490 \\
Gene--pathway fused & TCGA-BRCA & 3 & 0.256 & 0.287 & 0.000 & 0.566 \\
Gene--pathway fused & METABRIC & 3 & 0.258 & 0.206 & 0.000 & 0.440 \\
\bottomrule
\end{tabular}}
\end{table*}


The highest label-free view-ranking score was assigned to the pathway-fused representation, but it selected three clusters in both cohorts. In contrast, subtype-conditioned pathway signatures were highly reproducible across cohorts, demonstrating that biological signal was present even though the unsupervised selection criterion favored a coarser partition.

\subsection{Teacher--student ablation}
\begin{table*}[t]
\centering
\caption{Compact teacher--student ablation summary over ten development seeds. Balanced and unbalanced transport were numerically indistinguishable; the no-teacher and no-rare variants often improved agreement metrics.}
\label{tab:v52-appendix}
\resizebox{\textwidth}{!}{%
\begin{tabular}{llcccc}
\toprule
Variant & Cohort & NMI & ARI & Macro-F1 & Rare recall \\
\midrule
Full teacher--student & TCGA-BRCA & 0.268 $\pm$ 0.007 & 0.215 $\pm$ 0.006 & 0.398 & 0.175 \\
Full teacher--student & METABRIC & 0.279 $\pm$ 0.006 & 0.234 $\pm$ 0.008 & 0.402 & 0.078 \\
No teacher & TCGA-BRCA & 0.282 $\pm$ 0.004 & 0.254 $\pm$ 0.005 & 0.361 & 0.089 \\
No teacher & METABRIC & 0.275 $\pm$ 0.006 & 0.244 $\pm$ 0.007 & 0.386 & 0.059 \\
Balanced OT & TCGA-BRCA & 0.268 $\pm$ 0.007 & 0.215 $\pm$ 0.006 & 0.398 & 0.175 \\
Balanced OT & METABRIC & 0.279 $\pm$ 0.006 & 0.234 $\pm$ 0.007 & 0.402 & 0.079 \\
No OT & TCGA-BRCA & 0.263 $\pm$ 0.005 & 0.206 $\pm$ 0.006 & 0.393 & 0.181 \\
No OT & METABRIC & 0.278 $\pm$ 0.004 & 0.230 $\pm$ 0.008 & 0.393 & 0.066 \\
No rare weighting & TCGA-BRCA & 0.280 $\pm$ 0.010 & 0.233 $\pm$ 0.010 & 0.407 & 0.169 \\
No rare weighting & METABRIC & 0.276 $\pm$ 0.008 & 0.240 $\pm$ 0.009 & 0.418 & 0.088 \\
\bottomrule
\end{tabular}}
\end{table*}


The multi-view teacher was less stable and less confident than the single-view teacher, particularly in METABRIC. The full teacher--student configuration did not surpass the no-teacher variant, and the transport-balance and rare-weighting ablations did not support the corresponding mechanism claims.
